% Options for packages loaded elsewhere
% Options for packages loaded elsewhere
\PassOptionsToPackage{unicode}{hyperref}
\PassOptionsToPackage{hyphens}{url}
%
\documentclass[
  ignorenonframetext,
]{beamer}
\newif\ifbibliography
\usepackage{pgfpages}
\setbeamertemplate{caption}[numbered]
\setbeamertemplate{caption label separator}{: }
\setbeamercolor{caption name}{fg=normal text.fg}
\beamertemplatenavigationsymbolsempty
% remove section numbering
\setbeamertemplate{part page}{
  \centering
  \begin{beamercolorbox}[sep=16pt,center]{part title}
    \usebeamerfont{part title}\insertpart\par
  \end{beamercolorbox}
}
\setbeamertemplate{section page}{
  \centering
  \begin{beamercolorbox}[sep=12pt,center]{section title}
    \usebeamerfont{section title}\insertsection\par
  \end{beamercolorbox}
}
\setbeamertemplate{subsection page}{
  \centering
  \begin{beamercolorbox}[sep=8pt,center]{subsection title}
    \usebeamerfont{subsection title}\insertsubsection\par
  \end{beamercolorbox}
}
% Prevent slide breaks in the middle of a paragraph
\widowpenalties 1 10000
\raggedbottom
\AtBeginPart{
  \frame{\partpage}
}
\AtBeginSection{
  \ifbibliography
  \else
    \frame{\sectionpage}
  \fi
}
\AtBeginSubsection{
  \frame{\subsectionpage}
}
\usepackage{iftex}
\ifPDFTeX
  \usepackage[T1]{fontenc}
  \usepackage[utf8]{inputenc}
  \usepackage{textcomp} % provide euro and other symbols
\else % if luatex or xetex
  \usepackage{unicode-math} % this also loads fontspec
  \defaultfontfeatures{Scale=MatchLowercase}
  \defaultfontfeatures[\rmfamily]{Ligatures=TeX,Scale=1}
\fi
\usepackage{lmodern}

\usetheme[]{metropolis}
\ifPDFTeX\else
  % xetex/luatex font selection
  \setmathfont[]{Latin Modern Math}
\fi
% Use upquote if available, for straight quotes in verbatim environments
\IfFileExists{upquote.sty}{\usepackage{upquote}}{}
\IfFileExists{microtype.sty}{% use microtype if available
  \usepackage[]{microtype}
  \UseMicrotypeSet[protrusion]{basicmath} % disable protrusion for tt fonts
}{}
\makeatletter
\@ifundefined{KOMAClassName}{% if non-KOMA class
  \IfFileExists{parskip.sty}{%
    \usepackage{parskip}
  }{% else
    \setlength{\parindent}{0pt}
    \setlength{\parskip}{6pt plus 2pt minus 1pt}}
}{% if KOMA class
  \KOMAoptions{parskip=half}}
\makeatother


\usepackage{longtable,booktabs,array}
\usepackage{calc} % for calculating minipage widths
\usepackage{caption}
% Make caption package work with longtable
\makeatletter
\def\fnum@table{\tablename~\thetable}
\makeatother
\usepackage{graphicx}
\makeatletter
\newsavebox\pandoc@box
\newcommand*\pandocbounded[1]{% scales image to fit in text height/width
  \sbox\pandoc@box{#1}%
  \Gscale@div\@tempa{\textheight}{\dimexpr\ht\pandoc@box+\dp\pandoc@box\relax}%
  \Gscale@div\@tempb{\linewidth}{\wd\pandoc@box}%
  \ifdim\@tempb\p@<\@tempa\p@\let\@tempa\@tempb\fi% select the smaller of both
  \ifdim\@tempa\p@<\p@\scalebox{\@tempa}{\usebox\pandoc@box}%
  \else\usebox{\pandoc@box}%
  \fi%
}
% Set default figure placement to htbp
\def\fps@figure{htbp}
\makeatother





\setlength{\emergencystretch}{3em} % prevent overfull lines

\providecommand{\tightlist}{%
  \setlength{\itemsep}{0pt}\setlength{\parskip}{0pt}}



 


\setbeamerfont{title}{size=\large}
\usepackage{setspace}
\makeatletter
\@ifpackageloaded{caption}{}{\usepackage{caption}}
\AtBeginDocument{%
\ifdefined\contentsname
  \renewcommand*\contentsname{Table of contents}
\else
  \newcommand\contentsname{Table of contents}
\fi
\ifdefined\listfigurename
  \renewcommand*\listfigurename{List of Figures}
\else
  \newcommand\listfigurename{List of Figures}
\fi
\ifdefined\listtablename
  \renewcommand*\listtablename{List of Tables}
\else
  \newcommand\listtablename{List of Tables}
\fi
\ifdefined\figurename
  \renewcommand*\figurename{Figure}
\else
  \newcommand\figurename{Figure}
\fi
\ifdefined\tablename
  \renewcommand*\tablename{Table}
\else
  \newcommand\tablename{Table}
\fi
}
\@ifpackageloaded{float}{}{\usepackage{float}}
\floatstyle{ruled}
\@ifundefined{c@chapter}{\newfloat{codelisting}{h}{lop}}{\newfloat{codelisting}{h}{lop}[chapter]}
\floatname{codelisting}{Listing}
\newcommand*\listoflistings{\listof{codelisting}{List of Listings}}
\makeatother
\makeatletter
\makeatother
\makeatletter
\@ifpackageloaded{caption}{}{\usepackage{caption}}
\@ifpackageloaded{subcaption}{}{\usepackage{subcaption}}
\makeatother

\usepackage{bookmark}
\IfFileExists{xurl.sty}{\usepackage{xurl}}{} % add URL line breaks if available
\urlstyle{same}
\hypersetup{
  pdfauthor={Giorgio Arcara},
  hidelinks,
  pdfcreator={LaTeX via pandoc}}


\author{Giorgio Arcara}
\date{}

\begin{document}


\begin{frame}
% TITLE SLIDES
\title{Metodi Statistici per la Neuropsicologia Forense\\ \vspace{1em} \emph{3. La scelta dei test per la propria cassetta degli attrezzi}}
\author{Giorgio Arcara,\\ Università di Padova \\ IRCCS San Camillo, Venezia}

\titlegraphic{

\vspace*{7cm}
\includegraphics[scale=0.15]{Figures/LogoSanCamilloIRCCS_Unipd_alpha.png}
\hfill \includegraphics[scale=0.15]{Figures/CC_license_3_0.png}
}
\maketitle
\end{frame}

\begin{frame}{La scelta dei test}
\phantomsection\label{la-scelta-dei-test}
Due momenti di scelta dei test:

\begin{enumerate}[<+->]
[1)]
\item
  \underline{Scelta dei test da inserire nella propria «cassetta degli attrezzi»}
\item
  Scelta dei test da utilizzare con un determinato paziente
\end{enumerate}
\end{frame}

\begin{frame}{La scelta dei test da inserire nella propria «cassetta
degli attrezzi»}
\phantomsection\label{la-scelta-dei-test-da-inserire-nella-propria-cassetta-degli-attrezzi}
\emph{Regole generali da seguire}

Test che abbiano caratteristiche psicometriche soddisfacenti(validità,
affidabilità, dati normativi adeguati,
\underline{vedi slides successive})

Dipende molto da fattori pratici: il tipo di realtà clinica influenza
quali sono i test da utilizzare potenzialmente.
\end{frame}

\begin{frame}{La scelta dei test da inserire nella propria «cassetta
degli attrezzi»}
\phantomsection\label{la-scelta-dei-test-da-inserire-nella-propria-cassetta-degli-attrezzi-1}
\emph{Come selezionare i test?}

\underline{Possibili fonti:}

\textbf{Kit del Neuropsicologo}: un libro prodotto dalla Società
Italiana di Neuropsicologia (per i Soci) che contiene riferimenti a
quasi tutti I test disponibili in Italiano (non è però aggiornatissimo
ed è difficile da trovare).

\textbf{Libro ``L'esame neuropsicologico dell'Adulto''} (a cura di
Bianchi A.): contiene un capitolo che inlcude una lista piuttosto ampia
di test neuropsicologici adattati in Italiano.

\textbf{Corso: Valutazione delle funzioni cognitive in ambito forense}:
Presso il vostro corso di laurea.
\end{frame}

\begin{frame}{La scelta dei test da inserire nella propria «cassetta
degli attrezzi»}
\phantomsection\label{la-scelta-dei-test-da-inserire-nella-propria-cassetta-degli-attrezzi-2}
\emph{Come selezionare i test?}

\underline{Possibili fonti:}

\textbf{Case editrici di test (Organizzazioni Speciali, Eriksson)}: sono
test a pagamento che garantiscono un certo standard di qualità.

\begin{center}
  \textbf{METODI CHE PREFERISCO}
\end{center}

\textbf{Docenti universitari, futuri colleghi ed esperti di un certo
settore/patologia}: è il modo migliore (se trovate qualcuno disponibile
a condividere le informazioni).

\textbf{Articoli scientifici/consensus specifici per patologia}: Cercate
(es. su scholar o pubmed, se ci sono linee guida o criteri e cercate se
I test sono validati/standardizzati in Italiano).
\end{frame}

\begin{frame}{La scelta dei test da inserire nella propria «cassetta
degli attrezzi»}
\phantomsection\label{la-scelta-dei-test-da-inserire-nella-propria-cassetta-degli-attrezzi-3}
\textbf{Articoli scientifici/consensus}

Avere l'articolo scientifico spesso non basta. Serve anche avere il
manuale di istruzioni e di somministrazione (e attribuzione dei
punteggi) e i materiali del test, eventuali tabelle dei dati normativi

Utilizzate con cautela test di cui non avete referenze ufficiali
(possono avere modifiche ``custom'', che alterano il test)

\underline{Contattate gli autori del test (la mail è presente nell’articolo) e chiedete se è disponibile.}
\end{frame}

\begin{frame}{La scelta dei test da inserire nella propria «cassetta
degli attrezzi»}
\phantomsection\label{la-scelta-dei-test-da-inserire-nella-propria-cassetta-degli-attrezzi-4}
\emph{Cosa evitare nella selezione dei test?}

\begin{itemize}[<*>] 
  \item Evitare test troppo datati (se i dati normativi sono troppo vecchi potrebbero non essere utilizzabili) –     \emph{approfondiremo questo}.
  \item Evitare test “fatti in casa”: un test passato da un collega senza dati normativi (magari tradotti da un test in inglese) può essere accattivante (e talvolta utile) ma è limitato nelle informazioni che può fornire. - \emph{approfondiremo questo}.
\end{itemize}
\end{frame}

\begin{frame}{I criteri per cui scegliere i test per valutazioni in
ambito forense}
\phantomsection\label{i-criteri-per-cui-scegliere-i-test-per-valutazioni-in-ambito-forense}
\begin{columns}
\column{0.7\textwidth}

\textbf{1) Test adeguati alle potenziali domande della valutazioni che dovrete fare}

\column{0.3\textwidth}

\begin{figure}
    \includegraphics[scale=0.04]{Figures/triangle.png}
\end{figure}

\end{columns}

\pause

\begin{columns}
\column{1\textwidth}

\textbf{2)Test con buone capacità psicometriche}
\begin{itemize}[<*>] 
  \item Test con buona validità
  \item Test con buona affidabilità
  \item Test con buone proprietà di altro tipo (es. Valori di cambiamento)
  \item Test con dati normativi adeguati.
\end{itemize}

\column{0\textwidth}
\end{columns}

\vspace{2em}

In generale non dimenticate il razionale: I test ci forniscono
informazioni e parte di queste informazioni derivano dalle
caratteristiche dei test.
\end{frame}




\end{document}
